
\mysection[BurntOrange]{\centering Integral}



\DEF{4.1.1 (1)}{ $[a, b] \subseteq \mathbb{R}, f:[a, b] \rightarrow \mathbb{R}^n$ stetig.
Dann $$\int_a^b f(t) d t:=\left(\int_a^b f_1(t) d t, \ldots, \int_a^b f_n(t) d t\right) \in \mathbb{R}^n$$


}

\mysubsection{\centering Wegintegrale}
\NOTE{MC}{
\begin{itemize}
\item Falls $\delta$ eine orientierte Umparametrisierung von $\gamma$ ist mit $\varphi$. So ist auch $\gamma$ eine oU für $\delta$ mit $\varphi^{-1}$
\item Falls $f$ und $g$ ein Potential besitzen, dann besitzt die Summe $f+g$ ebenfalls ein Potential.
\end{itemize}
}

\KRZ{Wegintegral}{
\begin{enumerate}[leftmargin = *]
\item Parametrisiere $\gamma$ mit \COL{Parametrisierung}
\item Testen ob Berechnung mit \COL{S 4.6.3 (Satz von Green)} einfacher ist.
\item Berechne $\vec{\gamma}^{\prime}(t)=\frac{d}{d t} \vec{\gamma}(t)$ (jede Komponente nach dem Parameter $t$ ableiten).
\item Benutze die Formel:
$$
\int_\gamma \vec{v} \cdot d \vec{s}=\int_a^b \vec{v}(\vec{\gamma}(t)) \cdot \vec{\gamma}^{\prime}(t) d t .
$$

\end{enumerate}
}

\KRZ{Parametrisierungen}{
\begin{itemize}[leftmargin = *]
\item Funktion:  $y=f(x)$ oder $x = g(y)$
\begin{itemize}
\item $x=t, y=f(t)$
\item $x=g(t), y=t$
\end{itemize}
\item Liniensegment: Von $p_0$ nach $p_1$ mit $p_i = \left(x_i, y_i, z_i\right)$ 
\begin{itemize}
\item $\vec{r}=(1-t) p_0 +t p_1$
\end{itemize}
\item Kreis: Def $x^{2c} + y^{2c}= r^{2c}$
\begin{itemize}
\item UZS: $x=r \cdot \cos^{\frac{1}{c}} t, y=-r \cdot \sin^{\frac{1}{c}} t$ 
\item GUZS: $x=r \cdot \cos^{\frac{1}{c}} t, y=r \cdot \sin^{\frac{1}{c}} t$
\end{itemize}
\item Elipse: Def $\frac{x^{2c}}{a^{2c}}+\frac{y^{2c}}{b^{2c}}=1$
\begin{itemize}
\item UZS: $x=a \cdot \cos^{\frac{1}{c}} t, y=-b \cdot \sin^{\frac{1}{c}} t$
\item GUZS: $x=a \cdot \cos^{\frac{1}{c}} t, y=b \cdot \sin^{\frac{1}{c}} t$
\end{itemize}

\end{itemize}

UZS = Uhrzeigersinn, \\
GUZS = Gegenuhrzeigersinn\\
$c$: $a^{\frac{1}{c}}=\sqrt[c]{a}$ 
}

\KRZ{Potential}{
Gegeben sei ein konservatives Vektorfeld $f:\mathbb{R}^3 \rightarrow \mathbb{R}^3$ mit $f=(f_1, f_2,f_3)^\T$
Gesucht: $g: \mathbb{R}^3 \rightarrow \mathbb{R}$ mit $\nabla g=f$.

\begin{enumerate}
\item Definiere $g(x,y,z) = \int f_1 dx + a(y,z)$
\item Löse nach $a(y,z)$ in $\frac{\partial g}{\partial_y}=f_2$ und integriere $ a(y,z) = \int a(y,z) dy + b(z)$
\item Löse nach $b(z)$ in $\frac{\partial g}{\partial_z}=f_3$ und integriere $ b(z) = \int b(z)$
\end{enumerate}
Falls nach \hl{dem} Potenzial gesucht wird ist die Lösung $g+C$ wenn nach \hl{einem} Potenzial gesucht wird: $g$ mit $C=0$
}

\KRZ{Konservativität}{
\begin{itemize}[leftmargin =*]
\item Kontraposition \COL{Prop 4.1.13}: $\mathcal{J}_f$ nicht symmetrisch $\Rightarrow$ $f$ nicht konservativ
\item Falls $\mathcal{J}_f$ symmetrisch und Gebiet Sternförmig, ist Feld konservativ. (\COL{S 4.1.17})
\item Konstante Felder sind immer konservativ

\end{itemize}
}

\NOTE{Länge eines Weges}{
Für einen Weg $\gamma$ ist die Länge des Weges
$s = \int_a^b \sqrt{1+\left(\gamma^{\prime}(t)\right)^2} d t$

}

\DEF{4.1.1 (2)}{Ein \COL{Weg} ist ein stetiges $\gamma:[a, b] \rightarrow \mathbb{R}^n$. Wir betrachten nur \COL{stückweise} $C^1$ Wege, d.h. es gilt $k \geq 1$ and $a=t_0<t_1<\cdots<t_k=b$ Sodass $\left.\gamma\right|_{ \left( t_{i-1},t_{i}\right) } C^1$ für alle $1 \leq i \leq h$. Somit gibt es $k-1$ "Knickpunkte". Wir sagen $\gamma$ ist ein Weg von $\gamma(a)$ mach $\gamma(b)$.}

\DEF{4.1.1 (3)}{Ist $\gamma:[a, b] \rightarrow \mathbb{R}^n$ in weg, $U \subseteq \mathbb{R}^n$ sodass $\operatorname{Bild}(\gamma) \subseteq u$ und $f: U \rightarrow \mathbb{R}^n$ stetig. Dann ist das Wegintegral von $f$ entlang $\gamma$, geschrieben $\int_\gamma f(s) \cdot d \vec{s}$, definiert als
 
 $$
 \int_a^b f(\gamma(t)) \cdot \gamma^{\prime}(t) d t
 $$
 
 Eine alternative Scheibweise ist das Skalarprodukt
 $$
 \left\langle f(\gamma(t)), \gamma^{\prime}(t)\right\rangle
 $$}
 
 \DEF{4.1.4}{Eine \COL{orientierte Umparametrisierung} eines Wegs $\gamma:[a, b] \rightarrow \mathbb{R}^n$ ist ein Weg $\sigma:[c, d] \rightarrow \mathbb{R}^u$ mit $\sigma=\gamma \circ \varphi$, wobei $\varphi:[c, d] \rightarrow[a, b]$ stetig ist, differenzierbar auf $(c, d)$, streng monoton wachsend mit $\varphi(c)=a$ and $\varphi(d)=b$ (insbesondere bijektiv).}

\PROP{4.1.5}{$\gamma$ ein Weg in $\mathbb{R}^n$ mit einer orientierten Umparametrisierung $\sigma$ von $\gamma$ 
Sei weiter das Bild $\gamma \subseteq u \subseteq \mathbb{R}^n$ und $f: U \rightarrow \mathbb{R}^n$ stetig. Dann: $$
\int_\gamma f(s) \cdot d \vec{s}=\int_\sigma f(s) \cdot d \vec{s}
$$}

\DEF{4.1.8}{ $u \subseteq \mathbb{R}^m, \quad f: u \rightarrow \mathbb{R}^n$ stetig. Falls für alle $x_1, x_2 \in U$ das Integral
 $$
 \int_\gamma f(s) \cdot d \vec{s}
 $$
 unabhängig von der Wahl des Weges $\gamma$ von $x_1$ mach $x_2$ mit Bild $(\gamma) \subseteq U$ ist, dann ist $f$ \COL{konservativ.}}
 
 \DEF{}{$U$ heisst \COL{wegzusammenhängend} falls für alle $x, y \in U$ in Weg $\gamma$ von $x$ mach $y$ existiert mit $\operatorname{Bild}(\gamma)\subseteq U$.}
 
 \SA{4.1.10}{$u \subseteq \mathbb{R}^n$ offen, $f: u \rightarrow \mathbb{R}^n$ konservativ. Dann gibt es $\sin g \in C^1(u, \mathbb{R})$ mit $f=\nabla g$. Ein solches $g$ heisst \COL{Potential} von $f$ Falls $U$ wegzusammenhängend ist, dann is $g$ eindeutig bis auf Addition einer Konstanten.}
 
 
 \PROP{4.1.13}{$U \subseteq \mathbb{R}^n$ offen, $f: u \rightarrow \mathbb{R}^n C^1$.
$f$ konservativ $\Rightarrow \frac{\partial f_i}{\partial x_j}=\frac{\partial f_j}{\partial x_i}$ für $\quad i, j=1, \ldots, n$. $\Leftrightarrow J_f(x)$ symmetrisch für alle $x \in U$}


 \PROP{4.1.17}{$U \subseteq \mathbb{R}^n$ offen, \COL{sternförmig} $f: u \rightarrow \mathbb{R}^n C^1$.
$ \frac{\partial f_i}{\partial x_j}=\frac{\partial f_j}{\partial x_i}$ für $\quad i, j=1, \ldots, n\Rightarrow f$ konservativ }

\SA{}{$U \subseteq \mathbb{R}^n$ offen und \hl{sternförmig}, $f: U \rightarrow \mathbb{R}^n C^1$ mit $\frac{\partial f_i}{\partial x_j}=\frac{\partial f_j}{\partial x_i}$ für alle $i, j=1, \ldots, n \Rightarrow f$ konservativ}

\SA{}{$u \subseteq \mathbb{R}^\mu, f: u \rightarrow \mathbb{R}^n$ Konservativ $\Longleftrightarrow \oint_\sigma f(s) \cdot d \vec{s}=0$ für alle geschlossenen Wege $\gamma:[a, b] \rightarrow U$.}

\DEF{}{$U \subseteq \mathbb{R}^n$ heisst \COL{sternförmig} falls es ein $x_0 \in U$ gibt sodass für alle $x \in C$ die Strecke zwischen $x_0$ und $x$ in A enthalten ist.}

\NOTE{}{Es gilt Konvex $\Rightarrow$ Sternförmig $\Rightarrow$ wegzusammenhängend , da man zuerst immer zu $x_{0}$ gehen kann und dann zu einem beliebigen Punkt in $U$. }

\DEF{4.1.20}{$U \subseteq \mathbb{R}^3$ offen, $f: U \rightarrow \mathbb{R}^3 C^1$. Dann ist die Rotation das $C^0$ vektorfeld $u \rightarrow \mathbb{R}^3$ gegeben als
 $$
 \operatorname{rot}(f)=\left(\begin{array}{l}
 \partial_y f_3-\partial_z f_2 \\
 \partial_z f_1-\partial_x f_3 \\
 \partial_x f_2-\partial_y f_1
 \end{array}\right)
 $$
 Es gilt $J_f(x)$ symmetrisch $\Leftrightarrow \operatorname{rot}(f)(x)=0$}

\NOTE{}{Falls Domäne sternförmig ist gilt: Feld ist konservativ$ \Leftrightarrow J_f(x)$ symmetrisch $\Leftrightarrow \operatorname{rot}(f)(x)=0$
Zusammenfassend gilt:
\begin{gather*}  

f \in C^{0}(U,\mathbb{R}^{n})\text{ konservativ mit }U \subseteq \mathbb{R}^{n}\\  
\Updownarrow\\  
f=\nabla g\text{ für }g\in C^{1}(U,\mathbb{R})\\  
\qquad\qquad\qquad\qquad \Downarrow \quad(\Uparrow \quad falls U \text{ sternförmig})\\  
J_{f}(x) \text{ symmetrisch}
\end{gather*}

 }



\mysubsection{\centering Riemann Integral}

\KRZ{Normalbereiche}{
Für Mengen $\Omega$ wobei

$$
\Omega=\left\{(x, y) \in \mathbb{R}^2 \mid a \leq x \leq b, f(x) \leq y \leq g(x)\right\}
$$
und mit stetigen Funktionen $f, g$ und sei $F \in C^0(\bar{\Omega})$. Dann gilt

$$
\int_{\Omega} F d \mu=\int_a^b \int_{f(x)}^{g(x)} F(x, y) d y d x .
$$

Das gleiche gilt analog für Bereiche auf $f(y) \leq x \leq g(y)$ mit
$$
\int_{\Omega} F d \mu=\int_c^d \int_{h(y)}^{k(y)} F(x, y) d x d y .
$$
Ist $F$ stetig kann man die Reihenfolge der Integrale auch beliebig vertauschen.\\
Beispiele sind
\begin{itemize}[leftmargin = *]
\item Kreis: $\int_{0}^{\sqrt{r}}\int_{-\sqrt{r-x^2}}^{\sqrt{r-x^2}}f(x,y)dydx$
\end{itemize}
}

\DEF{}{
Für $U \subseteq \mathbb{R}^n$ kompakt, $f: U \rightarrow \mathbb{R}$ stetig kann man das (Riemannintegral von $f$ auf $u \int_u f\left(x_1, \ldots, x_n\right) d x$ definieren, mit folgenden Eigenschaften:

\begin{enumerate}[leftmargin = *]
\item \COL{Kompabilität} $n=1, U=[a, b] \Rightarrow \int_{[a, b]} f(x) d x=\int_a^b f(x) d x$
\item \COL{Linearität}  $\int_u a f(x)+b g(x) d x=a \int_u f(x) d x+b \int_u g(x) d x$ für $f, g: U \rightarrow \mathbb{R}$ stetig, $a, b \in \mathbb{R}$.
\item \COL{Positivität} $f \leqslant g \Rightarrow \int_u f(x) d x \leq \int_u g(x) d x$ und $f \geqslant 0, V \leq U$ kompakt $\Rightarrow \int_V f(x) d x \leq \int_U f(x) d x$
\item \COL{Dreiecksungleichung} $\left|\int_u f(x)+g(x) d x\right| \leq \int_u|f(x)| d x+\int_u|g(x)| d x$ und $\int_u f(x) d x \leqslant \int_u|f(x)| d x$	
\item \COL{Volumen} $\int_u 1 d x=\text { Volumen von } U$
\item \COL{Satz von Fubini} $n_1, n_2 \geqslant 1$ mit $n_1+n_2=n$. Sei
	$$
	\begin{aligned}
	& U  \subseteq \mathbb{R}^{n} = \mathbb{R}^{n_1}  \times \mathbb{R}^{n_2}\\  
	& V=\left\{x_1 \in \mathbb{R}^{n_1} \mid \exists x_2:\left(x_1, x_2\right) \in U\right\} \subseteq \mathbb{R}^{n_1} \\
& W\left(x_1\right)=\left\{x_2 \in \mathbb{R}^{n_2} \mid\left(x_1, x_2\right) \in U\right\} \subseteq \mathbb{R}^{n_2}
	\end{aligned}
	$$
	
	Falls $ g: V \longrightarrow \mathbb{R}, \quad g\left(x_1\right)=\int_{W\left(x_1\right)} f\left(x_1, x_2\right) d x_2 $ stetig ist, dann:
	$$
	\begin{aligned}
\int_u f(x) d x & =\int_V g\left(x_1\right) d x_1 \\
& =\iint_{W\left(x_1\right)} f\left(x_1, x_2\right) d x_2 d x_1
\end{aligned}
	$$
	
\item \COL{Addivität bezüglich Integrationsbereichs}: $U_1, U_2 \subseteq \mathbb{R}^n$ kompakt, $f: U_1 \cup U_2 \rightarrow \mathbb{R}$ stetig
	$
\Rightarrow \int_{u_1 \cup u_2} f(x) d x=\int_{u_1} f(x) d x+\int_{u_2} f(x) d x-\int_{u_1 \cap u_2} f(x) d x
$


\end{enumerate}



}

\NOTE{Volumina}{
\begin{itemize}[leftmargin=*]
\item Kreis: $A = \pi r^2$, wobei $r =$ Radius
\item Dreieck: $A = \frac{1}{2} b h$, wobei $b=$ Basis und $h=$ Höhe des Dreiecks
\item Kugel: $V = \frac{4}{3} \pi r^3$, wobei $r=$ Radius
\item Pyramide: $V = \frac{1}{3} s^2 h$, wobei $s =$ Seitenlänge der quadratischen Basis und $h= $ Höhe der Pyramide
\item Kegel: $V = \frac{1}{3} \pi r^2 h$, wobei $r=$ Radius der Basis und $h= $ Höhe des Kegels ist.
\item Ellipsoid: $V = \frac{4}{3} \pi a b c$, wobei $a$, $b$ und $c$ die Halbachsen des Ellipsoids sind.
\end{itemize}



}

\DEF{4.2.3}{
\begin{enumerate}
\item  Für $1 \leqslant m \leqslant n$ ist eine parametrisierte $m$-Menge in $\mathbb{R}^n$ eine stetige Funktion
 $$
 f:\left[a_1, b_1\right] \times \cdots \times\left[a_m, b_m\right] \rightarrow \mathbb{R}^n
 $$
die $C^1$ auf $\left(a_1, b_1\right) \times \cdots \times\left(a_m, b_m\right)$ ist.
 
 \item $B \subseteq \mathbb{R}^n$ heisst  \COL{vernachlässigbar} falls
 $B \subseteq B_i / d\left(f_1\right) \cup \cdots \cup$ Bild $\left(f_k\right) \quad$ für $m_i$-Mengen $f_i$ mit $m_i<n$.
\end{enumerate}
}

\NOTE{Def 4.2.3 Informell}{ Wir können die ganze Menge mit
einer endlichen Menge an Rechtecken beliebiger
Grösse überdecken.)}

\PROP{4.2.5}{ Ist $U \subseteq \mathbb{R}^n$ kompakt und vernachlässigbar, dann $\int_u f(x) d x=0$ für alle $f \in C^0(u, \mathbb{R})$.}

\mysubsection{\centering Uneigentliche Integrale}

\NOTE{Limes in $\mathbb{R}$}{Falls Integrationsbereich $(- \infty, \infty)$ aufteilen in $(-\infty ,c]$ und $[c, \infty)$ und separat integrieren. Falls $+\infty - \infty$ nicht definiert. (Bsp: $x$)}

\DEF{}{ Wir definieren zwei Arten von uneigentlichen Integrale
\begin{enumerate}[leftmargin = *]
\item $f:[a, \infty) \times I \rightarrow \mathbb{R}$ 
\item $f: \mathbb{R}^2 \rightarrow \mathbb{R}$
\end{enumerate}

Es gilt
\begin{enumerate}[leftmargin = *]

\item $[c,d]$ ist ein kompaktes Intervall, $a \in \mathbb{R}$, und $f:[a, \infty) \times [c,d] \rightarrow \mathbb{R}$ stetig. Dann

\begin{align*}
&\int_{[a, \infty) \times [c,d]} f(x, y) d x d y:\\ &=\lim _{b \rightarrow \infty} \int_{[a, b] \times [c,d]} f(x, y) d x d y\\
 &\overset{\underset{\mathrm{(1)}}{}}{=}\lim _{b \rightarrow \infty} \int_a^b \int_c^d f(x, y) d y d x\\ &=\int_a^{\infty} \int_c^d f(x, y) d y d x
 \\
 & \overset{\underset{\mathrm{(2)}}{}}{=}
  \int_c^d  \underbrace{\int_a^{\infty} f(x, y) d x}_{:=g(y)} d y
\end{align*}

Das Integral konvergiert falls das Limit existiert ($\neq \pm \infty$)


\item Für $f: \mathbb{R}^2 \rightarrow[0, \infty)$ (Funktionen die nur positiv sind) stetig ist, dann
\begin{align*}
&\int_{\mathbb{R}^2} f(x, y) d x d y\\&:=\lim _{R \rightarrow \infty} \int_{B_R(0)} f(x, y) d x d y \\&= \lim _{S \rightarrow \infty} \int_{[-S, S]^2} f(x, y) d x d y\\
 &\overset{\underset{\mathrm{(1)}}{}}{=} \lim_{S \to \infty} \int_{-S}^{S}\int_{-S}^{S} f(x,y) dy dx\\
 &\overset{\underset{\mathrm{(2)}}{}}{=} \int_{-\infty}^{\infty}\underbrace{\int_{-\infty}^{\infty} f(x,y) dy}_{:=g(x)} dx
\end{align*}
wobei $B_R(0)$ der Kreis um $0$ mit Radius $R$ ist. Das Gleiche erhält man wenn man Anstatt eine Scheibe ein Quadrat nimmt. 
\begin{enumerate}
\item[(1)] Fubini
\item[(2)] falls $g(q)$ für alle $q \in \mathbb{R}$ konvergiert und stetig ist.

\end{enumerate}
\end{enumerate}
}



\NOTE{Negative Funktionen}{ Für Funktionen die Negativ sind: Zwei Funktionen definieren 

\begin{align*}
f_{+}(x)=\begin{cases}
f(x) & \text{ falls } f(x)\geq 0 \\
0 & \text{sonst}
\end{cases}\\
 f_{-}(x)=\begin{cases}
-f(x) & \text{ falls } f(x)\leq 0 \\
0 & \text{sonst}
\end{cases}
\end{align*}
Das Ergebnis ist dann
$$\int f(x) dx = \int f_{+} dx -  \int f_{-} dx$$

}


\mysubsection{\centering Transformationsformel}


\KRZ{Transformation}{

\begin{itemize}
\item Definiere die Transformation 
$  \begin{bmatrix}
\phi_u(x,y) \\ \phi_v(x,y)
\end{bmatrix} =\begin{bmatrix}
u \\ v
\end{bmatrix} $
\item Benutze die Transformation um die Ränder des Integrals anzupassen
\item Berechne die Inverse der Transformation
$  \begin{bmatrix}
\phi_x(u,v) \\ \phi_y(u,v)
\end{bmatrix} =\begin{bmatrix}
x \\ y
\end{bmatrix} $
\item Berechne $Det = \operatorname{Det}\mathcal{}J_{\phi(u,v)^{-1}}=\operatorname{Det} \begin{bmatrix}
\partial x \phi_x & \partial y \phi_x \\
\partial x \phi_y & \partial y \phi_y \\
\end{bmatrix}$
\item Benutze die Transformationsformel mit den angepassten Rändern. \hl{$Det$ nicht vergessen}
\end{itemize} 



} 
\KRZ{Inverse Transformation}{
Gegeben sei die Transformationsfunktion
$\phi(x,y)= ...= (\phi_u(x,y) , \phi_v(x,y)) = (u,v)^T$
Wir bekommen die Gleichungen
$\phi_u(x,y)=u (I)$ und $\phi_v(x,y)=v (II)$.



Löse  $(I)$ nach $y$ auf und Setze in $(II)$ ein $\Rightarrow$ inverse $x=..$

Löse  $(II)$ nach $x$ auf und Setze in $(I)$ ein $\Rightarrow$ inverse $y=..$

}

\KRZ{Determinante}{
Für $2\times 2$ Matrizen gilt
$$\operatorname{det}\left|\begin{array}{ll}a & c \\ b & d\end{array}\right|=a d-b c$$
Für $3 \times 3$ Matrizen nach Zeile/Reihe entwickeln und abweichslungsweise addieren/subtrahieren. (Man started bei $+$)

}

\KRZ{Standard Transformationen}{
Schreibe das Integral um entsprechend der Transformationsformel.

\begin{itemize}[leftmargin=*]
\item Polar: $x^2 + y^2 = r^2 \Rightarrow x= r \cos(\theta), y = r \sin(\theta)$ mit $0\leq r < \infty, 0\leq\theta < 2 \pi$ und $\det = r$
\item Elliptisch: $\frac{x^2}{a^2}+\frac{y^2}{b^2}=1 \Rightarrow x= ar \cos(\theta), y = br \sin(\theta)$ mit $0\leq r <\infty, 0\leq\theta < 2 \pi$ und $\det =abr$	
\item Zylinder: $x^2 + y^2 = r^2 \Rightarrow  x= r \cos(\theta), y = r \sin(\theta), z=z$ mit $0\leq r < \infty,-\infty <z < \infty, 0\leq\theta < 2 \pi$ und $\det = r$
\item Kugel: $x^2 + y^2 + z^2 = r^2 \Rightarrow  x=r\sin(\theta)\cos(\varphi),y=r\sin(\theta)\sin(\varphi), z = r \cos(\theta)$ mit $0\leq r < \infty, 0\leq \theta < \pi, 0 \leq  \varphi < 2\pi$ und $\det = r^2 \sin(\theta)$
\end{itemize}

\\
\hl{Bei Radius $r^2$ Wurzel nicht vergessen}
Falls Bedingungen gestellt werden, Ränder entsprechen anpassen. (z.B. bei $\theta$)
Mann auch $-\pi \leq \theta < \pi$ anstatt $0\leq\theta < 2 \pi$ verwenden.
}

\NOTE{Winkel Anpassen}{
Falls spezifische Quadranten:
\begin{enumerate}[leftmargin = *]
\item  $x \geq 0$ (Obere Hälfte der Scheibe) : $0 \leq \theta\leq \pi$
\item   $y \geq 0$ (Linke Hälfte der Scheibe) : $\frac{\pi}{2} \leq \theta \frac{3\pi}{2}$
\item $x \geq 0,y \geq0$ : $0 \leq \theta\leq \frac{\pi}{2}$
\item $x \leq0,y \geq0$ : $\frac{\pi}{2} \leq \theta\leq \pi$
\item $x\leq0,y\leq0$ : $\pi \leq \theta\leq \frac{3\pi}{2}$
\item  $x\geq0,y\leq0$ : $\frac{3\pi}{2} \leq \theta\leq 2\pi$
\end{enumerate}
Andere Möglichkeiten:
\begin{enumerate}[leftmargin = *]
\item $x \geq y \geq0$ : $0 \leq \theta\leq \frac{\pi}{4}$
\item $y \geq x \geq0$ : $\frac{\pi}{4} \leq \theta\leq \frac{\pi}{2}$
\end{enumerate}
}


\KRZ{Komplexere Volumina}{
Sei $D=\left\{ ... \right\}$ eine beliebige Menge mit Bedingungen die alle von einer Koordinate $x_i$ abhängen. Dann
\begin{enumerate}[leftmargin=*]
\item Querschnittsfläche berechnen durch Umformen $W(x_i)$
\item Integrieren durch das Interval der Koordinate $x_i$
\end{enumerate}}

\DEF{}{$$
 \int_{\bar{U}} f(\varphi(x))\left|\operatorname{det} f_{\varphi}(x)\right| d x=\int_{\bar{V}} f(y) d y
$$
wobei:
\begin{itemize}

\item $\bar{U}$ kompakt, $\bar{U}=U \cup B$ für $U$ offen, $B$ vernachlässigbar
\item $\bar{V}$ kompakt, $\bar{V}=V \cup C$ für $V$ offen, $C$ vernachlässigbar
 \item $f: \bar{V} \rightarrow \mathbb{R}$ stetig}
\item $\varphi: \bar{U} \rightarrow \bar{V}$ stetig, und $C^1$ auf $U$
\item $\varphi(U)=V$ und $\varphi: U \to V$ bijektiv
\item $\left|\operatorname{det} f_{\varphi}(x)\right|$ lässt sich auf $\bar{U}$ stetig fortsetzen. (Funktion testen ob stetig auf $\bar{U}$)


\end{itemize}

\NOTE{Beispiele Transformation}{
$$
\int_{[0,a] \times[0,b]} f(x, y) d x d y=ab \int_0^1 \int_0^1 f(a x, b y) d x d y
$$

}

\mysubsection{Schwerpunktberechnung}

\DEF{}{Der \COL{Schwerpunkt} $\bar{x}\in \mathbb{R}^{n}$ einer kompakten Menge $U \subseteq \mathbb{R}^{n}$ ist gegeben durch
 $$\bar{x_{i}}=\frac{1}{\operatorname{Vol}(U)} \int_{U}x_{i}dx$$}
 
 
\NOTE{}{
Für einen Kegel mit beliebiger Grundfläche (z.B. Pyramide) liegt der Schwerpunkt immer auf $\frac{1}{4}$ von dem Schwerpunkt der Grundfläche zur Spitze.  

} 


\mysubsection{Satz von Green}	

In diesem Kapitel gilt: $\gamma$ stückweise $C^1$ und $\gamma ^{\prime}\neq \textbf{0}$ für alle $t \in[a, b]$ wo $\gamma^{\prime}(t)$ existiert.

\KRZ{Fläche mit Satz von Green}{
Gegeben sei eine Menge $C$ beschränkt durch $C_{\mathrm{pw}}^1$-Rand $\partial C$.
\begin{enumerate}[leftmargin =*]
\item Parametrisiere den Rand mit der Kurve $\gamma:[a, b] \rightarrow \mathbb{R}^2$: \COL{RC Parametrisierungen})
\item Wähle $v$ so dass $\operatorname{rot}(f)=1$: $\vec{v}=(0,x)$ oder $\vec{v}=(-y,0)$
\item Wende Satz von Green an:
$$
A = \int_{\gamma=\partial C} \vec{v} \cdot d \vec{s}=\int_a^b \vec{v}(\vec{\gamma}(t)) \cdot \vec{\gamma}^{\prime}(t) d t
$$

\end{enumerate}



}

\KRZ{Umgekehrter Satz von Green}{
Gegeben sei ein Linienintegral über einer Funktion $F$, für die $\left(\frac{\partial F_2}{\partial x}-\frac{\partial F_1}{\partial y}\right)$ simpel ist. Die Form des Linienintegrals ist idealerweise ein rechteck.

Berechne $\int_a^b \int_c^d \frac{\partial F_2}{\partial x}-\frac{\partial F_1}{\partial y} dx dy$, wobei das Rechteck die Eckpunkten $(a,c)$ und $(b,d)$ hat.



}

\DEF{4.6.1}{Eine einfache (= kreuzungsfreie) geschlossene Kurve in $\mathbb{R}^n$ ist ein Weg $\gamma:[a, b] \rightarrow \mathbb{R}^n$, sodass für $s, t \in[a, b]$ mit $s<t$ gilt: $\gamma(s)=\gamma(t) \Leftrightarrow s=a$ und $t=b$.}

\SA{Jordanischer Kurvensatz}{
Für $n=2\mathbb{R}^2 \backslash \operatorname{im}(\gamma)$ (wie definiert in 4.6.1) hat zwei Komponenten, eine beschränkte und eine umbeschränkte.
}

\SA{4.6.3 (Satz von Green)}{

$X \subseteq \mathbb{R}^2$ kompakt, mit Rand $\partial X$ gleich der disjunkten Vereinigung einfach geschlossener Jordankurven $\gamma_{1, \ldots,} \gamma_k$.
Ausserden liege $X$ stets links von $\gamma_i$.
$f: U \rightarrow \mathbb{R}^2$ ein $C^1$-Vektorfeld für in offenes $U$ mit $x \subseteq U$. Dann:

$$
\int_X \underbrace{\frac{\partial f_2}{\partial x}-\frac{\partial f_1}{\partial y}}_{\text{2-dim. Rotation}} d x d y=\sum_{i=1}^k \oint_{\gamma_i} f \cdot d \vec{s}
$$

}

\NOTE{}{
Mit \COL{S 4.6.3} kann eine beliebige Funktion $g$ integrieren. Dafür wählt man nur $f$, sodass
$$
\operatorname{rot}(f) = g
$$ 

}

\mysubsection{Lösungsmethoden}

\KRZ{Partielle Integration}{
Benutze Formel:
$$
I = \int_a^b u(x) v^{\prime}(x) d x=[u(x) v(x)]_a^b-\int_a^b v(x) u^{\prime}(x) d x .
$$

Für mehrere Iterationen benutze

\begin{center}

\begin{tabular}{l|ll}
    & $D$ & $I$ \\ \hline
$\textcolor[HTML]{1ebee6}{+}$   & $\textcolor[HTML]{1ebee6}{u}	$   & $v ^{\prime}$   \\
$\textcolor[HTML]{00e4c1}{-}$   &  $\textcolor[HTML]{00e4c1}{u^{\prime}}$ &  $\textcolor[HTML]{1ebee6}{v}$ \\ 
$\textcolor[HTML]{acfa70}{+}$   & $\textcolor[HTML]{acfa70}{u^{\prime\prime}}$  &   $\textcolor[HTML]{00e4c1}{\int v}$ \\ 
$-$   & $\vdots$ &  $\textcolor[HTML]{acfa70}{\vdots}$ \\ 
\end{tabular}
\end{center}

Das Resultat ist dann $\textcolor[HTML]{1ebee6}{uv} \textcolor[HTML]{00e4c1}{- u^{\prime} \int v} \ldots $.

Man führt die Tabelle solange bis:
\begin{itemize}[leftmargin = *]
\item eine Zeile integrierbar/konstant wird. 
\item eine Zeile $-I$ ergibt. (Dann umformen zu $2I = ...$ und lösen)
\end{itemize}




}

\KRZ{Integration durch Substitution}{
\begin{enumerate}[leftmargin = *]
\item Wähle substitution $\textcolor[HTML]{1ebee6}{u= g(x)}$
\item Bestimme $\textcolor[HTML]{1ebee6}{du = g(x) dx }$ und $\textcolor[HTML]{1ebee6}{dx = \frac{du}{g(x)}}$
\item Ersetze alle vorkommenden $x$-Terme
\end{enumerate}

\begin{align*}
\int_{a}^{b} f(x) dx &= \int_{\textcolor[HTML]{1ebee6}{g(a)}}^{\textcolor[HTML]{1ebee6}{g(b)}} f(x)[g(x)\mapsto u] \cdot  \textcolor[HTML]{1ebee6}{\frac{du}{g'(x)}} 
\\ &= \int_{g(a)}^{g(b)} f(u) du
\end{align*}

}
	
 

\end{flushleft}


